\section{A Toy Model Illustrating Non-Injective Projection}
  \label{app:toy-model}

  We present a simple toy model intended to illustrate how Bell-type correlations may
  arise from non-injective projection without invoking nonlocal dynamics, hidden
  variables, or superdeterministic assumptions.
  The model is deliberately minimal and does not aim to reproduce the full structure
  of quantum theory.
  Its purpose is solely to make explicit the logical mechanism underlying the failure
  of probabilistic factorization.

  \subsection{Underlying configuration space}
    \label{subsec:underlying-configuration-space}

    Let $\Omega$ denote a finite set of underlying configurations.
    Each configuration $\omega \in \Omega$ represents a complete specification of a
    relational structure.
    No spatiotemporal interpretation is assumed at this level, and no notion of
    subsystem decomposition is imposed.

    We consider a situation in which two measurement regions, denoted $A$ and $B$, are
    identified only at the level of observable descriptions.
    At the level of $\Omega$, the configurations need not decompose into independent
    components associated with $A$ and $B$.

  \subsection{Non-injective projection}
    \label{subsec:non-injective-projection}

    Observable outcomes are defined through a projection
    \[
      \Pi : \Omega \rightarrow O_A \times O_B ,
    \]
    where $O_A$ and $O_B$ denote the sets of possible outcomes for measurements performed
    in regions $A$ and $B$, respectively.

    The projection $\Pi$ is assumed to be non-injective.
    Distinct underlying configurations may therefore correspond to the same pair of
    observable outcomes $(a,b)$.
    Observable states are thus identified with equivalence classes
    \[
      [\omega] = \{ \omega' \in \Omega \mid \Pi(\omega') = \Pi(\omega) \}.
    \]

  \subsection{Emergence of correlations}
    \label{subsec:emergence-of-correlations}

    Let $\mu(\omega)$ be a probability distribution defined on $\Omega$.
    The joint probability distribution accessible at the observable level is given by
    \[
      P(a,b) = \sum_{\omega \in \Pi^{-1}(a,b)} \mu(\omega).
    \]

    Because the preimage $\Pi^{-1}(a,b)$ does not generically decompose into independent
    subsets associated with $A$ and $B$, the joint probability $P(a,b)$ does not admit a
    factorized representation of the form $P(a)P(b)$ or
    $P(a \mid \lambda) P(b \mid \lambda)$ for any effective variable $\lambda$.

    The resulting correlations arise not from any interaction or information exchange
    between $A$ and $B$, but from the fact that both outcomes are constrained by the same
    underlying relational configuration.
    From the perspective of the effective description, the two outcomes appear correlated
    even when the corresponding measurements are spacelike separated.

  \subsection{Free choice and absence of superdeterminism}
    \label{subsec:free-choice-and-absence-of-superdeterminism}

    Measurement settings at $A$ and $B$ may be freely and independently chosen within the
    effective description.
    The probability distribution $\mu(\omega)$ is assumed to be independent of these
    choices.
    No correlation between measurement settings and underlying configurations is required
    or assumed.

    The violation of statistical independence arises solely from the non-injective nature
    of the projection $\Pi$.
    Because observable outcomes do not provide access to the full underlying
    configuration, conditioning on observable variables does not isolate independent
    degrees of freedom.
    The failure of probabilistic factorization is therefore structural rather than
    conspiratorial.

  \subsection{Interpretation}
    \label{subsec:interpretation}

    This toy model illustrates how two observables may appear correlated because they
    represent different projections of a single underlying structure.
    In this sense, Bell-type correlations may be understood as consistency constraints
    imposed by a non-factorizable effective description, rather than as signatures of
    nonlocal causation.

    While highly simplified, the model captures the essential logical mechanism discussed
    in the main text.
    More elaborate realizations may be constructed by enriching the structure of
    $\Omega$ or the projection $\Pi$, but the origin of the correlations remains the same:
    the non-injective mapping between underlying configurations and observable outcomes.
